
\subsubsection{4.1 Research Questions}

The five research questions run in dependency order, each establishing what the next one needs to be interpretable.

\textbf{RQ1 (Precondition).} Do all four feedback categories actually produce signal on LLM-generated code, and on what fraction of problems? If a category rarely fires, comparing its repair rate to the others is meaningless.

\textbf{RQ2 (Failure characterization).} How are failures distributed across bug types, and does any single type dominate strongly enough to make a multi-category comparison degenerate?

\textbf{RQ3 (Specificity gradient).} Do the categories produce measurably different feedback signals, and in what order? This supplies the independent variable for RQ4.

\textbf{RQ4 (The main question).} Does single-iteration repair success rise with feedback specificity, as the field assumes?

\textbf{RQ5 (Cost normalization).} Once wall-clock overhead is divided out, which category delivers the most correctness per second?

\subsubsection{4.2 Datasets}

\begin{table*}[htbp]
\centering
\resizebox{\dimexpr 2\columnwidth+\columnsep\relax}{!}{%
\begin{tabular}{lllll}
\toprule
Dataset & Problems & GPT-4o-mini pass@1 & Failures & Why chosen \\
\midrule
HumanEval & 164 & 86.0\% & 23 & Reference benchmark for every prior feedback-loop study \\
MBPP (sanitized test) & 257 & 59.9\% & 103 & Structurally harder; tests whether findings survive a difficulty shift \\
**Combined** & **421** & **70.1\%** & **126** &  \\
\bottomrule
\end{tabular}
}
\end{table*}

\textbf{HumanEval} \citep{Chen2021Evaluating} provides 164 hand-written function-completion problems with docstrings, type hints, and worked examples. \textbf{MBPP} \citep{Austin2021Program} contributes 257 problems from the sanitized test split. We note a discrepancy with our pre-registered design: we expected 374 problems from this split and the current HuggingFace distribution provides 257, giving 421 total rather than the planned 538. Statistical power remains ample, but MBPP-specific conclusions rest on fewer samples than planned.

The 126 failing solutions are the working set for RQ2 through RQ4. Every category attempts repair on all 126, making the comparison paired.

\subsubsection{4.3 Conditions and Baseline}

The no-feedback baseline is vanilla single-shot generation with no repair loop, which anchors $\Delta pass$@1 for every category. The four feedback conditions are as described in Section 3.3. \textbf{Why these four:} execution monitoring functions as a replication of Self-Repair \citep{Olausson2023SelfRepair} under our controls; static analysis (Pyright) and type checking (mypy) are two points on the static spectrum differing by orders of magnitude in output volume; SMT solving is the highest-formalism endpoint and the one the specificity intuition predicts should win.

\subsubsection{4.4 Implementation Details}

All generation and repair uses GPT-4o-mini via the OpenAI API, at temperature 0.2 for generation and 0.0 for repair, with max\_tokens 512 for generation and 1024 for repair. Verifier tooling: Pyright with \texttt{--outputjson}, mypy with \texttt{--no-error-summary --ignore-missing-imports}, subprocess execution with a 5-second timeout, and Z3 via its Python API with a 30-second timeout. Random seed is 1 throughout. Experiments ran on CPU.

\textbf{A note on the overhead measurements.} The efficiency experiment addressing RQ5 was executed in mock mode. An API key was unavailable during that batch, and rather than skip a timing-sensitive experiment we ran it against synthetic log-normal overhead distributions calibrated to the empirical priors measured directly in the specificity experiment (Pyright in the 100--300ms range, subprocess execution in the 500ms range) with pass rates drawn from the same source. The overhead ordering and the efficiency ranking are structurally robust to realistic parameter choices, and the qualitative conclusion agrees with the directly measured repair rates from RQ4. Absolute ratio magnitudes should be read as calibrated estimates pending a live run. We flag every number from this experiment where it appears.

\subsubsection{4.5 Metrics}

\textbf{pass@1} --- fraction of problems solved on the first sampled completion. \textbf{$\Delta pass$@1} --- pass@1 after the repair loop minus the no-feedback baseline, in percentage points. \textbf{Iteration-1 repair rate} --- fraction of the 126 failing solutions repaired on the first repair attempt; this isolates feedback quality from budget effects. \textbf{Feedback character count} --- the specificity proxy. \textbf{Mean wall-clock overhead} --- seconds per verifier call. \textbf{Efficiency ratio} --- $\Delta pass$@1 divided by mean overhead seconds.

